\begin{papercard}{gilardi2023}
\cf{Stages} \stagetag{L} (zero-shot ChatGPT as annotator), \stagetag{I} (accuracy versus a trained-annotator gold standard; ``intercoder agreement'' as percent agreement between two runs or two coders). The PNAS paper that popularised LLM annotation. Four datasets (6,183 texts), five coding tasks, three annotator groups (trained research assistants, MTurk crowd-workers, ChatGPT), and two metrics. It is short (3 pages of results) and its design is simple, which makes it a useful baseline and a good object for critique.

\cf{Question and headline} Does zero-shot ChatGPT outperform crowd-workers on typical text-annotation tasks? Claims: accuracy exceeds MTurk by about 25 percentage points on average, ChatGPT's agreement exceeds that of crowd-workers and of trained annotators on all tasks, and cost per annotation is below \$0.003, roughly thirty times cheaper than MTurk \pc{gilardi2023}{1}{the per-annotation cost of ChatGPT is less than \$0.003}. Context: trained annotators give high-quality but expensive data, crowd-workers are cheaper but may be insufficient for complex tasks and non-English languages \pc{gilardi2023}{2}{Trained annotators tend to produce high-quality data, but involve significant costs}. The authors say it is the first systematic evaluation.

\cf{Data and tasks} Four samples: (i) 2,382 tweets from a corpus of 2.6 million on content moderation (Jan 2020--Apr 2021); (ii) 1,856 tweets by members of the US Congress (2017--2022) from 20 million; (iii) 1,606 newspaper articles on content moderation (LexisNexis, from 980k); (iv) a new January 2023 sample of 500 tweets (339 English) added ``to address the concern that ChatGPT might be relying on memorization'' \pc{gilardi2023}{5}{Sample size was determined by the number of texts needed to build a training set for a machine learning classifier}. Tasks: relevance (content moderation; political content), topic (six classes), stance on repealing Section 230 (three classes), general frames (problem/solution/neither) and policy frames (14 classes from the Media Frames Corpus). Not every task applies to every dataset; attempts at human annotation of stance, topic and policy frames for news articles failed \pc{gilardi2023}{8}{for news articles, our attempts at human annotation for stance, topic, and policy frames were not successful}.

\cf{Annotator groups} \emph{Trained annotators}: three political-science students coded independently, using the project's codebook. \emph{Crowd-workers}: MTurk, restricted to ``Masters'', approval rate above 90\,\% with at least 50 approved HITs, located in the US; no worker could annotate more than 20\,\% of the tweets for a task; each text annotated by two workers \pc{gilardi2023}{5}{we ensured that no worker could annotate more than 20\% of the tweets for a given task}. \emph{ChatGPT}: API model gpt-3.5-turbo, annotations 9--20 March and 27 April--4 May 2023; the \emph{same} instruction text as for humans, deliberately without ChatGPT-specific prompting, and each tweet submitted individually with a prompt of the type ``Here's the tweet I picked, please label it as [task instruction]'' \pc{gilardi2023}{5}{We intentionally avoided adding any ChatGPT-specific prompts to ensure comparability between ChatGPT and MTurk crowd-workers}; temperature 1 (default) and 0.2, two responses each, so four responses per text; a new chat session per tweet to avoid history effects \pc{gilardi2023}{6}{We created a new chat session for every tweet to ensure that the ChatGPT results are not influenced by the history of annotations}.

\cf{Metrics} \emph{Accuracy}: share of correct annotations with trained annotators as gold standard, ``considering only texts that both annotators agreed upon''; \emph{intercoder agreement}: percentage of instances for which two annotators in a group (or two ChatGPT runs) give the same class \pc{gilardi2023}{6}{intercoder agreement refers to the percentage of instances for which both annotators in a given group report the same class}. No chance-corrected coefficient and no confidence intervals are reported.

\cf{Results} ChatGPT outperforms MTurk in accuracy for most tasks (Figure 1). Accuracy on two-class relevance tasks ranges 59--83\,\% (e.g.\ 83\,\% US Congress tweets, 59\,\% 2023 content-moderation tweets; the 2023 relevance score suffered from misclassifying suspension tweets because of ``a lack of examples in the prompt'') \pc{gilardi2023}{4}{ChatGPT performed much better than MTurk in the second task but struggled with misclassifying tweets about specific user suspensions in the relevance task due to a lack of examples in the prompt}. Average intercoder agreement: MTurk 56\,\%, trained annotators 79\,\%, ChatGPT temperature 1: 91\,\%, temperature 0.2: 97\,\% \pc{gilardi2023}{4}{On average, intercoder agreement is about 56\% for MTurk, 79\% for trained annotators, 91\% for ChatGPT with temperature = 1, and 97\% for ChatGPT with temperature = 0.2}. Correlations across tasks: agreement and accuracy of ChatGPT $r=.36$; ChatGPT accuracy and trained-annotator agreement $r=.46$ (better on easier tasks); ChatGPT's \emph{advantage over MTurk} and trained-annotator agreement $r=-.37$ (larger advantage on harder tasks) \pc{gilardi2023}{4}{ChatGPT’s outperformance of MTurk is negatively correlated with the intercoder agreement of trained annotators}. The authors call the test ``hard'' because the categories were designed for their own research and several tasks have many classes.

\cf{Discussion and future work} Evidence is ``consistent across different types of texts and time periods'', suggesting ChatGPT may already be superior to crowd-annotation; open questions: other languages, few-shot learning, semi-automated labelling in which a model learns from human annotations, chain-of-thought, and comparison across LLMs \pc{gilardi2023}{4}{(i) performance across multiple languages; (ii) implementation of few-shot learning}. Replication materials on Harvard Dataverse \pc{gilardi2023}{6}{Replication materials are available at the Harvard Dataverse}.

\cf{Limits (card; several are acknowledged by later work such as Pangakis et al.\ and Törnberg)} (1) ChatGPT's ``intercoder agreement'' is the agreement of repeated runs of the \emph{same model}; it measures output consistency (and falls with temperature), not independent-coder reliability, so it cannot be compared with human inter-coder agreement; consistency can coexist with systematic error. (2) Percent agreement is not chance-corrected and is inflated by few classes and skewed prevalence; the paper reports no $\kappa$/$\alpha$. (3) The gold standard only keeps texts on which both trained coders agreed, which removes the hardest items and can inflate the apparent accuracy of every group and bias the human-vs.-LLM comparison; the comparison on those items is also not the performance on the full data. (4) Gold standard of three students for categories the authors themselves developed, one study, one topic (content moderation), English tweets/news; no formal power or uncertainty estimate. (5) Single model version (gpt-3.5-turbo, 2023) that is no longer available; the exact API output cannot be reproduced. (6) MTurk was configured with a no-extra-prompt, one-pass design and may be weaker than a well-managed crowd. (7) Only unambiguous, single-label tasks.

\cf{Use for this thesis} (1) The reference showing why the thesis' stage-L comparison must be designed with care: reporting only percent agreement of an LLM with itself, or with a gold standard that excludes disagreements, proves little. The thesis should compare LLM, coder~1 and coder~2 on the same units, report chance-corrected $\alpha$/$\kappa$ per code and keep contested items in. (2) The temperature observation (lower temperature = higher self-consistency, no loss of accuracy) supports fixing temperature low, but test it. (3) A reminder to pose the task to humans and model with the identical codebook text and to log date and model version. (4) Because it is the best-known positive result, cite it for motivation but pair it with cautionary work (Pangakis et al., Baumann et al., Törnberg) when making claims.
\end{papercard}
